%%%%%%%%%%%%%%%%%%%%%%%%%%%%%%%%%%%%%%%%%%%%%%%%%%%%%%%%%%%%%%%%%%%%%%%%%%%%%%
%% reference:
%% http://www-i6.informatik.rwth-aachen.de/~dreuw/latexbeamerposter.php
%%%%%%%%%%%%%%%%%%%%%%%%%%%%%%%%%%%%%%%%%%%%%%%%%%%%%%%%%%%%%%%%%%%%%%%%%%%%%%

\documentclass[final]{beamer}
\mode<presentation>
{
  %\usetheme{I6dv}
  \usetheme{l6dv-2}
  %\usetheme{l6dv-3}
%  \usetheme{Aachen}
%  \usetheme{Oldi6}
%  \usetheme{I6td}
  %\usetheme{I6dv}
%  \usetheme{I6pd}
%  \usetheme{I6pd2}
}


\usepackage{listings}
% this part is used to allow your readers to copy the code from a PDF viewer but without copying the line numbers.
\usepackage{accsupp}
\newcommand*{\noaccsupp}[1]{\BeginAccSupp{ActualText={}}#1\EndAccSupp{}}

\lstdefinestyle{shared}
{
    numbers=none,
    frame=single,
    framesep=\fboxsep,
    framerule=\fboxrule,
    rulecolor=\color{red},
    xleftmargin=\dimexpr\fboxsep+\fboxrule\relax,
    xrightmargin=\dimexpr\fboxsep+\fboxrule\relax,
    breaklines=true,
    tabsize=2,
    columns=flexible,
}

\lstdefinestyle{python}
{
    style=shared,
    language={Python},
    basicstyle=\tt\footnotesize,
    keywordstyle=\color{blue},
    commentstyle=\color[rgb]{0.13,0.54,0.13},
    backgroundcolor=\color{gray!2},
    morekeywords={
        Console,
        WriteLine,
        int,
    },
}

\lstnewenvironment{python}
{\lstset{style=python}}
{}



% additional settings
\usefonttheme[onlymath]{serif}
\setbeamerfont{itemize}{size=\normalsize}
\setbeamerfont{itemize/enumerate body}{size=\normalsize}
\setbeamerfont{itemize/enumerate subbody}{size=\small}
\setbeamertemplate{bibliography entry title}{}
\setbeamertemplate{bibliography entry location}{}
\setbeamertemplate{bibliography entry note}{}

\usepackage{natbib}

%\renewcommand\printleftlogo
%  {\begin{center}
%     \resizebox{\textwidth}{!}%
%       {\includegraphics{GMU_Logo.eps}}
%   \end{center}
%  }

% additional packages
\usepackage{times}
\usepackage{amsmath,amsthm, amssymb, latexsym,epsfig,graphicx}
\usepackage{mathtools}
\usepackage{exscale}
\usepackage{multirow}
%\usepackage{subfig}
\usepackage{color}

\usepackage{pgfplots}
\pgfplotsset{every axis/.append style={
                    xlabel={$x$},          % default put x on x-axis
                    ylabel={$y$},          % default put y on y-axis
                    label style={font=\tiny},
                    tick label style={font=\tiny}
            }}
\pgfplotsset{every tick label/.append style={font=\tiny}}
\pgfplotsset{every axis/.append style={font=\tiny}}
\pgfplotsset{compat=1.12}



\usepackage{subcaption}
\captionsetup{compatibility=false}
\usepackage{booktabs, array}
%\usepackage{tabularx}
%\usepackage{rotating} %sideways environment
\usepackage[english]{babel}
%\usepackage[latin1]{inputenc}
%\usepackage[orientation=landscape,size=custom,width=200,height=120,scale=1.9]{beamerposter}
\usepackage[size=custom, width=48,height=36,scale=1.4]{beamerposter}
\listfiles


\title{\LARGE Chebyshev Expansions: Background and Computing in the
C++/Python Library ChebTools}
\author[Bouck, Bell]{Lucas Bouck and Ian Bell}
\institute[George Mason University]{George Mason University, Fairfax, VA, U.S.A. and National Institute of Standards and Technology, Boulder, CO, U.S.A}
%\date[Aug. 31 , 2007]{Aug. 31 , 2007}

% abbreviations
\usepackage{xspace}
\makeatletter
%\DeclareRobustCommand\onedot{\futurelet\@let@token\@onedot}
%\def\@onedot{\ifx\@let@token.\else.\null\fi\xspace}
%\def\eg{{e.g}\onedot} \def\Eg{{E.g}\onedot}
%\def\ie{{i.e}\onedot} \def\Ie{{I.e}\onedot}
%\def\cf{{c.f}\onedot} \def\Cf{{C.f}\onedot}
%\def\etc{{etc}\onedot}
%\def\vs{{vs}\onedot}
%\def\wrt{w.r.t\onedot}
%\def\dof{d.o.f\onedot}
%\def\etal{{et al}\onedot}
\makeatother

%%%%%%%%%%%%%%%%%%%%%%%
% Logo for Poster
%%%%%%%%%%%%%%%%%%%%%
%\printleftlogo
%\leftlogo[1]{GMU_Logo.eps}  % defines logo to left of title (with scale factor)
%\rightlogo[0.7]{nsf_logo.jpg}  % same but on right

%%%%%%%%%%%%%%%%%%%%%%%%%%%%%%%%%%%%%%%%%%%%%%%%%%%%%%%%%%%%%%%%%%%%%%%%%%%%%%%%%%%%%%%%%%%%%%%%%%%%%%%%%%%%
%%%%%%%%%%%%%%%%%%%%%%%%%%%%%%%%%%%%%%%%%%%%%%%%%%%%%%%%%%%%%%%%%%%%%%%%%%%%%%%%%%%%%%%%%%%%%%%%%%%%%%%%%%%%
\begin{document}
%\addtobeamertemplate{headline}{} 
%{
%\begin{tikzpicture}[remember picture,overlay]
%\node [shift={(-5.5 cm,-4cm)}] at (current page.north east) {\includegraphics[height=6.5cm]{GMU_Logo}}; 
%\end{tikzpicture} 
%}
\vspace*{-10mm}
\begin{frame}[fragile]{}
  \begin{columns}[t]
    \hspace*{5mm}

    \begin{column}{.31\linewidth}

    %%%%%%%%%%%%%%%%%%%%%%%%%%%%%%%%%%%%%%%%%%%%%%%%%%%%%%%%%%%%%%%%%%%%%%%%%%%%%%%%%%%%%%%%%%%%%%%%%%%%%%%%%%%%
\begin{block}{Summary}
	This poster has two goals:
	\begin{enumerate}
	\item Give a brief background of Chebyshev interpolation
		\begin{itemize}
		\item Background on Polynomial Interpolation
		\item Motivation and useful properties Chebyshev Polynomials
		\end{itemize}
	\item Introduce and demonstrate the capabilities of {\tt ChebTools}
		\begin{itemize}
		\item C++/Python library for working with Chebyshev expansions
		\item In the early stages and are looking for help in growing the library
		\end{itemize}	
	\end{enumerate}
\end{block}
\vspace*{2mm}
\begin{block}{Polynomial Interpolation}
{\color{numhypRed} What is it}\\
\vspace{2mm}
	\begin{minipage}{.55\linewidth}
		\begin{itemize}
		\item Given points, $\{(x_i,y_i)\}_{i=0}^{n}$ 
		\item Want a degree $n$ polynomial, $p_n$
		\item $p_n(x_i) = y_i$ for all $i$
		\item Interpolate a function $f$ by setting $y_i = f(x_i)$
		\item Gives an approximation of a function with just addition and multiplication for operations
	\end{itemize}
	\end{minipage}%}
	\hfill%
	\begin{minipage}{.4\linewidth}
	\begin{figure}
	\includegraphics[width=\textwidth]{interpolation_ex.pdf}
	\end{figure}
	\end{minipage}%}
	\\
	\vspace{2mm}
{\color{numhypRed} Need to be Careful \cite[Thm.\ 6.1]{mason_2003}}\\
	\begin{minipage}{.55\linewidth}
	\begin{itemize}
		\item Want to interpolate $f(x) = \frac{1}{1+25x^2}$ on $[-1,1]$ 
		\item Use evenly spaced points, 
		\item Polynomials actually diverge as $n\to\infty$
	\end{itemize}
	\end{minipage}%}
	\hfill%
	\begin{minipage}{.4\linewidth}
	\begin{figure}
	\includegraphics[width=\textwidth]{runge_bad.pdf}
	\end{figure}
	\end{minipage}%}
\end{block}
\begin{block}{Motivation for Chebyshev Polynomials}
{\color{numhypRed} Polynomial Interpolation Error \cite[Thm.\ 3.1.1]{davis_1975}}\\
The error of polynomial interpolation for a function that is $C^{n+1}[a,b]$:
$$f(x)-p_n(x) = \frac{f^{n+1}(c)}{(n+1)!}\prod_{j=1}^{n+1}(x-x_j),$$														%%%%%%%%%%% cite %%%%%%%%%%%%%%%%%%%
{\color{numhypRed}  Motivating Fact \cite[Thm.\ 3.3.4]{davis_1975}}\\
For $[-1,1]$, the Chebyshev Roots
$$x_j=\cos\left(\frac{2j-1}{2(n+1)}\pi\right), 1\leq j\leq n+1$$
minimize the term $\max_{x\in[a,b]}\left|\prod_{j=1}^{n+1}(x-x_j)\right|$
\end{block}


\begin{block}{Chebyshev Polynomials}
\begin{minipage}{.5\textwidth}
{\color{numhypRed} Definition}
$$T_n(x)=\cos\left(n\arccos(x)\right),\, x\in[-1,1]$$
{\color{numhypRed} Recurrence Relation}
\begin{align*}
T_0(x) &= 1\quad T_1(x) = x\\
T_{n+1}(x) &= 2xT_n(x)-T_{n-1}(x)
\end{align*}
Recurrence allows for fast evaluation of Chebyshev polynomial expansions
\end{minipage}
\begin{minipage}{.45\textwidth}
\begin{center}
\begin{figure}
\includegraphics[width=.8\textwidth]{cheb_poly_plot.pdf}
\caption{First 5 Chebyshev Polynomials}
\end{figure}
\end{center}
\end{minipage}
\end{block}


\end{column}
%%% end of first column %%%%%%%%%%%%

%% 2/19/2012 start next column next time %%%%%%%%%%%
%%%%%%%%%%%%%%%%%%%%%%%%%%%%%%%%%%%%%%%%%%%%%%%%%%%%%%%%%%%%%%%%%%%%%%%%%%%%%%%%%%%%%%%%%%%%%%%%%%%%%%%%%%%
%%% begining of second column %%%%%%%%%%%%%%
    \begin{column}{.31\linewidth}
    %%%%%%%%%%%%%%%%%%%%%%%%%%%%%%%%%%%%%%%%%%%%%%%%%%%%%%%%%%%%%%%%%%%%%%%%%%%%%%%%%%%%%%%%%%%%%%%%%%%%%%%%%%%%%%%%%%%%%%%%%%%%%%%%%%%%%%%%%%%

%

\begin{block}{Useful Properties of Chebyshev Polynomials}
{\color{numhypRed} Interpolation Convergence \cite[Ch.\ 8]{trefethen_2012}}\\    %%%%%%%%%%% cite %%%%%%%%%%%%%%%%%%%
\begin{minipage}{.4\textwidth}
\begin{itemize}
\item Suppose we want to interpolate $f:[-1,1]\to \mathbb{R}$
\item $f$ is analytic $\implies$ geometric convergence 
\item Need a small degree expansion for machine precision in our example
\end{itemize}
\end{minipage}
\begin{minipage}{.55\textwidth}
\begin{figure}
\includegraphics[width=\textwidth]{sine_error.pdf}
\caption{Error of Chebyshev Interpolation on $\sin(x)$ on $[0,2\pi]$}
\end{figure}
\end{minipage}
{\color{numhypRed} Orthogonality}\\
Chebyshev polynomials are orthogonal on $[-1,1]$ with a weight
$$\int_{-1}^{1}\frac{T_n(x)T_m(x)}{\sqrt{1-x^2}} = 0 \quad \text{if $m\neq n$}$$
Discrete orthogonality relations allow for quick interpolation \cite[Sec.\ 6.3]{mason_2003}.
\end{block}

\begin{block}{{\tt ChebTools} Introduction}
\begin{itemize}
\item Written in C++11 with wrappers in Python using {\tt pybind11} \cite{pybind11}
\item Backend linear algebra with {\tt Eigen} \cite{eigen}						%%%%%%%%%%	cite		%%%%%%%%%%%
\end{itemize}
{\color{numhypRed} Chebyshev Expansion Object}
If our expansion looks like 
$$\sum_{k=0}^{n}a_kT_k(x), \quad \text{on $[-1, 1]$}$$ 
then the object looks like
$$ {\tt xmin} = -1, \; {\tt xmax} = 1, \; {\tt coef} = [a_0, a_1, \ldots, a_n]$$
{\color{numhypRed} Related Libraries}
\begin{itemize}
\item {\tt chebfun} \cite{driscoll_2014}: Extensive MATLAB library by Nick Trefethen's group
\item {\tt ApproxFun.jl} \cite{approxfun}: Julia library by Alex Townsend
\item {\tt pychebfun}\cite{pychebfun} and {\tt chebpy}\cite{chebpy} for Python users
\end{itemize}

\end{block}

\begin{block}{{\tt ChebTools} Current Basic Capabilities}
{\color{numhypRed} Generating Expansions}\\
\begin{python}
# arguments: (degree of expansion, function, xmin, xmax)
cheb_sin = ct.generate_Chebyshev_expansion(10, np.sin, 0, 2*np.pi)
\end{python}
{\color{numhypRed} Function Evaluation}\\
\begin{python}
# works with floats and numpy arrays
cheb_sin.y(3)
\end{python}
{\color{numhypRed} Arithmetic Operations}\\
\begin{python}
# addition by another expansion and a constant
# note that operators like +=, -, and -= are also supported
chebs_added = cheb_sin + cheb_cos + 1
# multiplication of two expansions and multiplication by a constant
chebs_sin2x = 2*cheb_sin *cheb_cos
\end{python}
{\color{numhypRed} Applying More Advanced Functions (from {\tt pychebfun} \cite{pychebfun})}\\
\begin{python}
# applying more complicated functions to a Chebyshev expansion object
sin_squared = cheb_sin.apply(lambda x: x**2)
\end{python}\end{block}
\end{column}
%% end of second column %%%%%%%%%%%%%%%%%%%%%%%%%%%%%%%%%%%%%%%%%%%%%%%%%%%%%%%%%%%%


%%%% begining of third column%%%%%%%%%%%%%%%%%%%%%%%%%%%%%%%%%%%%%%%%%%%%%%%%%%%%%%%%%%%%
  \begin{column}{.31\linewidth}
\begin{block}{Other Current Capabilities}
{\color{numhypRed} Root Finding, Derivatives, and Local Extrema}\\

\begin{python}
# generate expansion of cosine
long_cos = ct.generate_Chebyshev_expansion(25, np.cos, -.1, 4*np.pi+.1)

# finding the roots of the Chebyshev expansion
# the boolean argument being True means that we want just the roots in [xmin, xmax]
approx_roots = long_cos.real_roots(True)

# finding the local extrema of the Chebyshev expansion
approx_extrema = long_cos.deriv(1).real_roots(True)
\end{python}
%\begin{center}
%\begin{figure}
%\caption{Cosine with its approximate roots on $[-.1,4\pi+.1]$}
%\includegraphics[width = .6\textwidth]{cos_roots.pdf}
%\end{figure}
%\end{center}

	\begin{minipage}{.4\linewidth}
		\begin{itemize}
		\item The method {\tt real\_roots} utilizes a companion matrix method \cite{boyd_2013}
		\item Faster method in the C++ library that avoids the companion matrix
		\item Has been used in \cite{bell_2018-2}
		\end{itemize}
	\end{minipage}%}
	\hfill%
	\begin{minipage}{.55\linewidth}
	\begin{figure}
	\includegraphics[width = \textwidth]{cos_roots.pdf}
	\end{figure}
	\end{minipage}%}
\end{block}



\begin{block}{Future Work}
We are looking for help with development. We want to add additional capabilities to the library including:
\begin{itemize}
\item Integrating expansions
\item 2-D capabilities
\item Adaptive Interpolation
\item Parallelization
\end{itemize}
\end{block}

\begin{block}{More Info}
{\color{numhypRed} JOSS Paper}\\
Bell et al., (2018). ChebTools: C++11 (and Python) tools for working with Chebyshev expansions. {\it ]ournal of Open Source Software}, 3(22), 569. https://doi.org/10.21105/joss.00569\\
{\color{numhypRed} GitHub Link }\\
https://github.com/usnistgov/ChebTools\\
{\color{numhypRed} Jupyter Notebook Example}\\
{\small https://github.com/usnistgov/ChebTools/blob/master/BattlesTrefethen.ipynb}
\end{block}

\begin{block}{References and Acknowledgements}
\bibliographystyle{plain}
\scriptsize{\bibliography{scipy2018_poster}}
\bigskip
{\normalsize This work was supported by the NIST SURF program.}
\end{block}
\end{column}
  \end{columns}
\end{frame}

\end{document}


%%%%%%%%%%%%%%%%%%%%%%%%%%%%%%%%%%%%%%%%%%%%%%%%%%%%%%%%%%%%%%%%%%%%%%%%%%%%%%%%%%%%%%%%%%%%%%%%%%%%
%%% Local Variables:
%%% mode: latex
%%% TeX-PDF-mode: t
%%% End: 